\section{Conduite à tenir selon l’étudiante }

Différents scénarios sont envisageables mais selon moi, le scénario le plus éthique aurait été une prise en charge multidisciplinaire dès la naissance du patient et la découverte du déni de grossesse. Une prise en charge psychologique voire psychiatrique selon le contexte, me parait obligatoire pour prévenir toute maltraitance et trouble dans la relation mère-enfant ainsi que de désamorcer la culpabilité que la maman ressentira. Prolonger l’hospitalisation lors de la naissance jusqu’à observer un attachement entre l’enfant et la mère ou tout signe qui va en ce sens afin d’avoir l’assurance nécessaire pour permettre un retour à domicile me parais indispensable. La suite du plan thérapeutique comprendrait une mise en place des aides dans le quotidien de la famille et en plus, être attentif en tant que médecin au deuil que la maman a traversé lors des premières semaines de vie du patient. Je pense que cette prise en charge aurait pu prévenir le passage en réanimation. Sans même transformer trop les faits et idéaliser la psychothérapie, la surveillance extérieure de cette famille, avec un regard bienveillant et médical aurait permis de dépister la dénutrition plus précocement. Avec cette présence du monde médical, cela permet de gagner du temps et laisser l’opportunité d’aborder l’approche centrée sur le patient.\\ 

Apporter cependant une nuance à ce scénario utopique est primordial car la condition nécessaire pour que cela fonctionne est que les parents acceptent la prise en charge et la présence de l’équipe soignante dans leur quotidien. A nouveau, avec une approche centrée sur le patient, les soignants ne s’imposent pas aux patients donc je reste convaincue qu’avec une communication adaptée une relation thérapeutique peut être entretenue. Il est important de notifier également que les parents se sont présentés à plusieurs consultations médicales prévues ce qui montre une collaboration fragile mais présente. Sans doute, ceci est une porte d’entrée pour un travail ensemble dans l’intérêt du patient.\\

\todoH{FloraStyle}{red!30}{pe mettre ce lien en note bas de page parce que ça m'a inspiré pour trouver une conduite en plus à faire \url{https://www.yapaka.be/sites/yapaka.be/files/texte/15_06_19_le_deni_de_grossesse-a_mettre_en_ligne_0.pdf}}{Le professionnel de santé} a bien sûr son rôle à jouer dans la reconnaissance de la grossesse le plus tôt possible afin d’éviter un déni de grossesse total. Au-delà du diagnostic précoce qui permet de limiter les complications, il est simplement important que le médecin reconnaisse le phénomène et la déconstruction des croyances qui entourent le déni de grossesse. Éviter le tabou de celui-ci permet une meilleure prise en charge et éviter la dénutrition sévère par carence d’apport. Des protocoles standardisés pour encadrer la prise en charge du déni de grossesse permettent de limiter les erreurs individuelles de chaque médecin et d’uniformiser les connaissances. La rédaction de ces protocoles serait l’idéal et une des premières idées qui vient en tête. Toutefois, si je dois rester réaliste, je dois considérer la difficulté que cela représente de mener des études cliniques et pouvoir tirer des conclusions statistiquement significatives avec un échantillon aussi réduit et ainsi pouvoir élaborer des guidelines pertinentes.\\ 

Pour conclure ce cas clinique, je pense que la prise en charge réelle n’a pas été optimale et surtout aurait pu être différente à plusieurs niveaux. Cependant, je pense qu’il est toujours plus simple de réfléchir au cas clinique a posteriori, sans urgence et avec du recul sur l’évolution de la situation. La prise de décision est toujours influencée par de nombreux facteurs et parfois, certains arguments prennent plus de poids qu’ils ne devraient ce qui entraine des conséquences néfastes. Pour ce patient, la suite de son histoire a été tragique mais grâce à plusieurs garde-fous, l’enfant a survécu et l’espoir d’une amélioration de son développement est toujours possible.\\  
